\documentclass[11pt,a4paper]{article}

\usepackage[utf8]{inputenc}
\usepackage[T1]{fontenc}
\usepackage[margin=2.5cm]{geometry}
\usepackage{booktabs}
\usepackage{xcolor}
\usepackage{listings}
\usepackage{hyperref}
\usepackage{parskip}

\hypersetup{colorlinks=true, linkcolor=black, urlcolor=blue}

% --- code listing style -------------------------------------------------
\definecolor{codebg}{rgb}{0.96,0.96,0.96}
\definecolor{codecomment}{rgb}{0.40,0.45,0.40}
\definecolor{codekw}{rgb}{0.10,0.20,0.60}
\lstset{
  basicstyle=\ttfamily\small,
  backgroundcolor=\color{codebg},
  commentstyle=\color{codecomment}\itshape,
  keywordstyle=\color{codekw}\bfseries,
  breaklines=true,
  frame=single,
  rulecolor=\color{gray!40},
  showstringspaces=false,
  columns=fullflexible,
  xleftmargin=0.5em,
  literate={\^}{{\textasciicircum}}1
}

\title{\textbf{Branch Integration and Bioplastic Class Extension}\\
\large CNNIRPLASTICS --- FTIR Polymer Classification}
\author{Jose Maria Contreras Prada}
\date{2 June 2026}

\begin{document}
\maketitle

\begin{abstract}
This report documents two changes made to the \texttt{CNNIRPLASTICS} repository:
(1) the integration of the \texttt{origin/main} development line into the
\texttt{add-cosine-baseline-eval} branch, unifying the preprocessing pipeline
with the new random-forest analysis and reporting work; and (2) the extension of
the classification target set with two bioplastics, PLA and PHA, at the
data-builder and label-definition level. Both changes were committed and pushed
to the remote.
\end{abstract}

\tableofcontents
\bigskip

% =======================================================================
\section{Branch Merge}
% =======================================================================

\subsection{Motivation}

Prior to this work the branch \texttt{add-cosine-baseline-eval} and the
\texttt{origin/main} branch had diverged: neither was a superset of the other.
The feature branch carried the end-to-end pipeline and spectral preprocessing
work (\texttt{run\_pipeline.py}, \texttt{preprocess.py}, the cosine-similarity
baseline, and the real-world lab samples), whereas \texttt{main} carried a
reworked random-forest pipeline, single-spectrum test-time estimation, the
feature-importance analysis, and an expanded requirements section in the README.
Merging was required so that subsequent retraining used a single, unified code
base.

\subsection{Procedure}

A safety branch \texttt{backup/pre-merge-add-cosine} was created before the
merge. \texttt{origin/main} was then merged into \texttt{add-cosine-baseline-eval}.
The merge produced content conflicts in three files, which were resolved
manually as summarised in Table~\ref{tab:conflicts}.

\begin{table}[h]
\centering
\caption{Conflict resolution decisions.}
\label{tab:conflicts}
\small
\begin{tabular}{@{}p{0.27\linewidth} p{0.36\linewidth} p{0.30\linewidth}@{}}
\toprule
\textbf{File} & \textbf{Nature of conflict} & \textbf{Resolution} \\
\midrule
\texttt{models/cnn\_model\_draft.py} &
Old monolithic \texttt{load\_and\_preprocess} (main) vs.\ refactored
\texttt{create\_model} / \texttt{train\_model} / \texttt{test\_model} (branch) &
Kept branch: the pipeline depends on the refactored API \\
\addlinespace
\texttt{models/rf\_model.py} &
\texttt{preprocess}-module routing (branch) vs.\ inline min--max plus
feature-importance saving (main) &
Kept main: preserves the new feature-importance analysis; its min--max is
equivalent to \texttt{normalize="minmax"} \\
\addlinespace
\texttt{data/format\_data.py} &
Four cosmetic print-string differences (unit notation, whitespace) &
Kept branch: strings are already pure ASCII and stdout is forced to UTF-8 \\
\bottomrule
\end{tabular}
\end{table}

\subsection{Verification}

After resolution, the absence of conflict markers was confirmed, all core
Python modules compiled, and the public APIs imported cleanly: the dataset
loader, the preprocessing entry points (including SNV and AsLS/arPLS), and the
refactored CNN functions. The six commodity-plastic classes (HDPE, LDPE, PP,
PS, PVC, PET) were unchanged by the merge.

\subsection{Outcome}

The merge was recorded as commit \texttt{6cd97ea} and pushed to the remote.
Files from both development lines are present in the unified tree: the pipeline
and preprocessing work from the feature branch, and the random-forest,
single-test, and requirements work from \texttt{main}. The pre-existing SNV and
AsLS/arPLS implementations in \texttt{preprocess.py} were not modified and remain
wired into the pipeline configuration.

% =======================================================================
\section{Bioplastic Class Extension (PLA and PHA)}
% =======================================================================

\subsection{Background}

The OpenSpecy spectral data committed to the repository contains only the six
commodity plastics. A search of the 8\,390 metadata identity records returned
zero matches for PLA- or PHA-related terms. This is by design: the dataset
builder script downloads the raw OpenSpecy library, filters it down to the six
target classes, and writes only the matched spectra to disk. The two bioplastics
were therefore absent from every committed artefact.

\subsection{Keyword Groups Added to the Dataset Builder}

Two new keyword groups were added to the \texttt{target\_polymers} list in
\texttt{data/data\_processing/openspecy\_dataset\_builder.R}. Each polymer is
matched against the lowercased \texttt{spectrum\_identity} field using
word-boundary regular expressions, so the new groups follow the same matching
convention as the existing six.

\begin{lstlisting}[language=R]
PLA = c(
  "pla",
  "polylactic acid",
  "poly\\(lactic acid\\)",
  "polylactide",
  "lactic acid",
  "poly-l-lactic acid",
  "plla"
),

PHA = c(
  "pha",
  "phb",
  "phbv",
  "polyhydroxyalkanoate",
  "polyhydroxyalkanoates",
  "polyhydroxybutyrate",
  "poly\\(hydroxybutyrate\\)",
  "poly-3-hydroxybutyrate",
  "polyhydroxybutyrate-co-valerate"
)
\end{lstlisting}

The PHA group intentionally aggregates the PHB and PHBV variants, which are
chemically members of the polyhydroxyalkanoate family, into a single class.

\subsection{Label Definition Update}

The two classes were registered in the canonical label list in
\texttt{data/format\_data.py}. Because the number of output classes is derived
from the length of this list throughout the code base, this single edit resizes
both the CNN softmax layer and the random-forest output space automatically.

\begin{lstlisting}[language=Python]
POLYMER_CLASSES = ["HDPE", "LDPE", "PP", "PS", "PVC", "PET", "PLA", "PHA"]
\end{lstlisting}

This yields eight classes, with integer indices \texttt{PLA = 6} and
\texttt{PHA = 7}.

\subsection{Current Status and Caveat}

The change is code-level scaffolding only. The two new classes remain
\emph{empty} until the OpenSpecy dataset is regenerated by re-running the builder
script, which requires an R environment and a re-download of the raw OpenSpecy
library. Until then no training spectrum carries a PLA or PHA label. The pipeline
remains runnable in this state---the two classes simply contribute no samples---
but a retrain should not be relied upon for bioplastic discrimination until the
data is regenerated. In addition, while PLA is expected to be well represented in
the OpenSpecy library, PHA/PHB/PHBV may be sparse; the class counts printed by
the builder should be inspected before the class is trusted, and a second data
source may be required.

\subsection{Outcome}

The keyword groups and the label update were recorded as commit \texttt{c20b143}
and pushed to the remote.

% =======================================================================
\section{Summary of Commits}
% =======================================================================

\begin{table}[h]
\centering
\caption{Commits produced by this work, on branch
\texttt{add-cosine-baseline-eval}.}
\label{tab:commits}
\small
\begin{tabular}{@{}l l p{0.50\linewidth}@{}}
\toprule
\textbf{Hash} & \textbf{Type} & \textbf{Description} \\
\midrule
\texttt{6cd97ea} & Merge & Integrate \texttt{origin/main} into the feature branch;
resolve three conflicts \\
\texttt{c20b143} & Feature & Add PLA and PHA keyword groups and label definitions
(scaffolding) \\
\bottomrule
\end{tabular}
\end{table}

\section{Next Steps}

\begin{enumerate}
  \item Regenerate the OpenSpecy dataset by running the extended builder script
  in an R environment; inspect the per-class spectrum counts.
  \item Clear the cached processed array (\texttt{output/\_cache\_processed.npz})
  so the new eight-class data is reprocessed from scratch.
  \item Retrain the CNN and random-forest models on the unified pipeline and
  evaluate per-class F1, with attention to the two new bioplastic classes.
  \item If PHA coverage is insufficient, identify a supplementary FTIR source
  for that class.
\end{enumerate}

\end{document}
